\section{AMC 10 10B 2002}
        \begin{enumerate}
        	\item (\textit{AMC 10 10B 2002}) The ratio $\frac{2^{2001}\cdot3^{2003}}{6^{2002}}$ is:


 $\mathrm{(A) \ } \frac{1}{6}\qquad \mathrm{(B) \ } \frac{1}{3}\qquad \mathrm{(C) \ } \frac{1}{2}\qquad \mathrm{(D) \ } \frac{2}{3}\qquad \mathrm{(E) \ } \frac{3}{2}$


 
 \textbf{Answer: }E



\textbf{Solution 1}

$\frac{2^{2001}\cdot3^{2003}}{6^{2002}}=\frac{6^{2001}\cdot 3^2}{6^{2002}}=\frac{9}{6}=\frac{3}{2}$ or $\mathrm{ (E) \ }$

 


 



\textbf{Solution 2}

$\frac{2^{2001}\cdot3^{2003}}{6^{2002}}=\frac{2^{2001}\cdot 2\cdot 3^{2002}\cdot 3}{6^{2002}\cdot 2}=\frac{2^{2002} \cdot 3^{2002} \cdot 3}{6^{2002}\cdot 2}=\frac{6^{2002}\cdot 3}{6^{2002}\cdot 2}=\frac{3}{2}$ or $\mathrm{ (E) \ }$~by mathwiz0

 



\textbf{Solution 3}

$\frac{2^{2001}\cdot3^{2003}}{6^{2002}}=\frac{2^{2001}\cdot3^{2003}}{2^{2002}\cdot3^{2002}}=\frac{3}{2}$

 


	\item (\textit{AMC 10 10B 2002}) For the nonzero numbers $a, b,$ and $c,$ define

 \[(a,b,c)=\frac{abc}{a+b+c}\]
Find $(2,4,6)$.


 $\mathrm{(A) \ } 1\qquad \mathrm{(B) \ } 2\qquad \mathrm{(C) \ } 4\qquad \mathrm{(D) \ } 6\qquad \mathrm{(E) \ } 24$


 
 \textbf{Answer: }C



\textbf{Solution}

$\frac{2\cdot 4\cdot 6}{2+4+6}=\frac{48}{12}=4\Longrightarrow\mathrm{ (C) \ }$

 


	\item (\textit{AMC 10 10B 2002}) The arithmetic mean of the nine numbers in the set $\{9,99,999,9999,\ldots,999999999\}$ is a $9$-digit number $M$, all of whose digits are distinct. The number $M$ does not contain the digit


 $\mathrm{(A) \ } 0\qquad \mathrm{(B) \ } 2\qquad \mathrm{(C) \ } 4\qquad \mathrm{(D) \ } 6\qquad \mathrm{(E) \ } 8$


 
 \textbf{Answer: }A



\textbf{Solution 1}

We wish to find $\frac{9+99+\cdots +999999999}{9}$, or $\frac{9(1+11+111+\cdots +111111111)}{9}=123456789$. This doesn't have the digit 0, so the answer is $\boxed{\mathrm{(A)}\ 0}$

 



\textbf{Solution 2}

Notice that the final number is guaranteed to have the digits $\{1, 3, 5, 7, 9\}$ and that each of these digits can be paired with an even number adding up to 9. $\boxed{\mathrm{(A)}\ 0}$ can be taken out, with the other digits fulfilling divisibility by 9.

 



\textbf{Solution 3}

The arithmetic mean is $\frac{(10^1-1)+(10^2-1)+\ldots+(10^9-1)}{9}=\frac{1111111101}{9}=123456789$. So select $\boxed{\mathrm{A}}$.~hastapasta

 


	\item (\textit{AMC 10 10B 2002}) What is the value of


 $(3x-2)(4x+1)-(3x-2)4x+1$


 



 when $x=4$?


 $\mathrm{(A) \ } 0\qquad \mathrm{(B) \ } 1\qquad \mathrm{(C) \ } 10\qquad \mathrm{(D) \ } 11\qquad \mathrm{(E) \ } 12$


 
 \textbf{Answer: }D



\textbf{Solution 1}

By the distributive property, 

 
 \[(3x-2)[(4x+1)-4x] + 1 = 3x-2 + 1 = 3x-1 = 3(4) - 1 = \boxed{\mathrm{(D)}\ 11}\].

 



\textbf{Solution 2}

Inputting 4 into $x$ in the original equation,

 
 \[[3(4)-2][4(4)+1]-[3(4)-2][4(4)]+1 = 170-160+1 = \boxed{\mathrm{(D)}\ 11}\]

 


	\item (\textit{AMC 10 10B 2002}) Circles of radius $2$ and $3$ are externally tangent and are circumscribed by a third circle, as shown in the figure. Find the area of the shaded region.


 
                    \begin{figure}[H]
                        \centering
                        \includegraphics[width=0.5\textwidth]{Problem/AMC_10_10B_2002_5_Problem_Figure_1.png}
                    \end{figure}
                    

 $\mathrm{(A) \ } 3\pi\qquad \mathrm{(B) \ } 4\pi\qquad \mathrm{(C) \ } 6\pi\qquad \mathrm{(D) \ } 9\pi\qquad \mathrm{(E) \ } 12\pi$


 
 \textbf{Answer: }E



\textbf{Solution 1}

A line going through the centers of the two smaller circles also goes through the diameter. The length of this line within the circle is $3+3+2+2=10.$ Because this is the length of the larger circle's diameter, the length of its radius is $5.$

 The area of the large circle is $25\pi$, and the area of the two smaller circles is $9\pi + 4\pi = 13\pi.$ To find the area of the shaded region, subtract the area of the two smaller circles from the area of the large circle. $\longrightarrow 25\pi - 13\pi = \boxed{\mathrm{(E) \ } 12\pi}$

 



\textbf{Solution 2}

The diameter of the larger circle is $10$ thus making the radius $5$ and the area $25\pi$. The smaller circle have radii of $2$ and $3$, giving us areas of $4\pi$ and $9\pi$ respectively. So $4\pi$ + $9\pi$ = $13\pi$ of the $25\pi$ is covered up meaning only $25\pi - 13\pi$ = $\boxed{\mathrm{(E) \ } 12\pi}$ of area is shaded.

 


	\item (\textit{AMC 10 10B 2002}) For how many positive integers $n$ is $n^2-3n+2$ a prime number?


 $\mathrm{(A) \ } \text{none}\qquad \mathrm{(B) \ } \text{one}\qquad \mathrm{(C) \ } \text{two}\qquad \mathrm{(D) \ } \text{more than two, but finitely many}\qquad \mathrm{(E) \ } \text{infinitely many}$


 
 \textbf{Answer: }B



\textbf{Solution 1}

Factoring, we get $n^2 - 3n + 2 = (n-2)(n-1)$. Either $n-1$ or $n-2$ is odd, and the other is even.  Their product must yield an even number.  The only prime that is even is $2$, which is when $n$ is $3$ or $0$. Since $0$ is not a positive number, the answer is $\boxed{\mathrm{(B)}\ \text{one}}$.

 



\textbf{Solution 2}

Considering parity, we see that $n^2 - 3n + 2$ is always even.  The only even prime is $2$, and so $n^2-3n=0$ whence $n=3\Rightarrow\boxed{\mathrm{(B)}\ \text{one}}$.

 


	\item (\textit{AMC 10 10B 2002}) Let $n$ be a positive integer such that $\frac{1}{2}+\frac{1}{3}+\frac{1}{7}+\frac{1}{n}$ is an integer. Which of the following statements is \textbf{not} true?


 $\mathrm{(A) \ } 2\text{ divides }n\qquad \mathrm{(B) \ } 3\text{ divides }n\qquad \mathrm{(C) \ } 6\text{ divides }n\qquad \mathrm{(D) \ } 7\text{ divides }n\qquad \mathrm{(E) \ } n>84$


 
 \textbf{Answer: }E



\textbf{Solution 1}

Since $\frac 12 + \frac 13 + \frac 17  = \frac {41}{42}$, $0 < \lim_{n \rightarrow \infty} \left(\frac{41}{42} + \frac{1}{n}\right) < \frac {41}{42} + \frac 1n < \frac{41}{42} + \frac 11 < 2$

 From which it follows that $\frac{41}{42} + \frac 1n = 1$ and $n = 42$. The only answer choice that is not true is $\boxed{\mathrm{(E)}\ n>84}$.

 



\textbf{Solution 2 (no limits)}

Since $\frac 12 + \frac 13 + \frac 17  = \frac {41}{42}$, it is very clear that $n=42$ makes the expression an integer*. Because $n$ is a positive integer, $\frac{1}{n}$ must be less than or equal to $1$. Thus, the only integer the expression can take is $1$, which means the only value for $n$ is $42$. Thus $\boxed{\mathrm{(E)}\ n>84}$

 



\textbf{Solution 2.1 (faster ending to solution 2)}

Once you find $n=42$, you can skip the rest of Solution 2 by noting that $n=42=2\times3\times7$, which satisfies A, B, C, D, but not E. Thus $\boxed{\mathrm{(E)}\ n>84}$ must be the answer by PoE (Process of Elimination).

 



\textbf{Solution 3(similiar to solution2)}

Cross multiplying and adding the fraction we get the fraction to be equal to $\frac {41n + 42}{42n}$, This value has to be an integer. This implies,

 $42n|(41n+42)$.

 => $42|(41n + 42)$

 \begin{verbatim}
but $42|42$, hence $42|41n$  but $42$ does not divides $41$,  $42|n$                                                                                                    -(1)
\end{verbatim}

=> $n|(41n+42)$

 \begin{verbatim}
but $n|41n$,  $n|42$                                                                                                    -(2)
\end{verbatim}

from (1) and (2) we get that n=42.Comparing this with the options, we see that option E is the incorrect statement and hence E is the answer.

 


	\item (\textit{AMC 10 10B 2002}) Suppose July of year $N$ has five Mondays. Which of the following must occur five times in the August of year $N$? (Note: Both months have $31$ days.)


 $\textbf{(A)}\ \text{Monday} \qquad \textbf{(B)}\ \text{Tuesday} \qquad \textbf{(C)}\ \text{Wednesday} \qquad \textbf{(D)}\ \text{Thursday} \qquad \textbf{(E)}\ \text{Friday}$


 
 \textbf{Answer: }D



\textbf{Solution 1}

If there are five Mondays, there are only three possibilities for their dates: $(1,8,15,22,29)$, $(2,9,16,23,30)$, and $(3,10,17,24,31)$. 

 In the first case August starts on a Thursday, and there are five Thursdays, Fridays, and Saturdays in August. 

 In the second case August starts on a Wednesday, and there are five Wednesdays, Thursdays, and Fridays in August.

 In the third case August starts on a Tuesday, and there are five Tuesdays, Wednesdays, and Thursdays in August.

 The only day of the week that is guaranteed to appear five times is therefore $\boxed{\textrm{(D)}\ \text{Thursday}}$.

 



\textbf{Solution 2 (visualization)}

A diagram can be constructed to visualize the possibilities for the dates. First, a 4x7 rectangle representing 28 days can be constructed with three extra cells on the bottom to represent the three extra days. 

 Then, the rest of the steps in Solution 1 can be followed.

 For example, if the first of the bottom three cells is a Monday, it is clear that the following two days are Tuesday and Wednesday. This means that the new bottom three cells (as both months have 31 days) of August will be Thursday, Friday, and Saturday. This same process can be repeated so that the second cell is represented as Monday and so forth; then, the only day of the week that appears when all three cases are considered is $\boxed{\textrm{(D)}\ \text{Thursday}}$.

 


	\item (\textit{AMC 10 10B 2002}) Using the letters $A$, $M$, $O$, $S$, and $U$, we can form five-letter "words". If these "words" are arranged in alphabetical order, then the "word" $USAMO$ occupies position


 $\mathrm{(A) \ } 112\qquad \mathrm{(B) \ } 113\qquad \mathrm{(C) \ } 114\qquad \mathrm{(D) \ } 115\qquad \mathrm{(E) \ } 116$


 
 \textbf{Answer: }D



\textbf{Solution 1}

There are $4!\cdot 4$ "words" beginning with each of the first four letters alphabetically. From there, there are $3!\cdot 3$ with $U$ as the first letter and each of the first three letters alphabetically. After that, the next "word" is $USAMO$, hence our answer is $4\cdot 4!+3\cdot 3!+1=\boxed{115\Rightarrow\text{(D)}}$.

 



\textbf{Solution 2}

Let $A = 1$, $M = 2$, $O = 3$, $S = 4$, and $U = 5$. The ways the numbers can be arranged is $5$ factorial which is $120$, with $12345$ in spot $1$ and $54321$ in spot $120$. Then counting backwards, $54321, 54312, 54231, 54213, 54132, 54123$, so the answer is $\boxed{115\Rightarrow\text{(D)}}$

 


	\item (\textit{AMC 10 10B 2002}) Suppose that $a$ and $b$ are nonzero real numbers, and that the equation $x^2+ax+b=0$ has solutions $a$ and $b$. What is the pair $(a,b)$?


 $\mathrm{(A) \ } (-2,1)\qquad \mathrm{(B) \ } (-1,2)\qquad \mathrm{(C) \ } (1,-2)\qquad \mathrm{(D) \ } (2,-1)\qquad \mathrm{(E) \ } (4,4)$


 
 \textbf{Answer: }C



\textbf{Solution 1}

Since $(x-a)(x-b) = x^2 - (a+b)x + ab = x^2 + ax + b = 0$, it follows by comparing \href{https://artofproblemsolving.com/wiki/index.php?title=Coefficient}{coefficients} that $-a - b = a$ and that $ab = b$. Since $b$ is nonzero, $a = 1$, and $-1 - b = 1 \Longrightarrow b = -2$. Thus $(a,b) = \boxed{\mathrm{(C)}\ (1,-2)}$.

 



\textbf{Solution 2}

Another method is to use \href{https://artofproblemsolving.com/wiki/index.php?title=Vieta\%27s\_formulas}{Vieta's formulas}. The sum of the solutions to this polynomial is equal to the opposite of the $x$ coefficient, since the leading coefficient is 1; in other words, $a + b = -a$ and the product of the solutions is equal to the constant term (i.e, $a*b = b$). Since $b$ is nonzero, it follows that $a = 1$ and therefore (from the first equation), $b = -2a = -2$. Hence, $(a,b) = \boxed{\mathrm{(C)}\ (1,-2)}$

 



\textbf{Solution 3}

Since $a$ and $b$ are solutions to the given equation, we can write the two equations $a^2 + a^2 + b = 2a^2 + b = 0,$ and $b^2 + ab + b = 0.$ From the first equation, we get that $b = -2a^2,$ and substituting this in our second equation, we get that $4a^4 - 2a^3 - 2a^2 = 0,$ and solving this gives us the solutions $a = -\frac{1}{2},$ $a = 0$ and $a = 1.$ We discard the first two solutions, as the first one doesnt show up in the answer choices and we are given that $a$ is nonzero. The only valid solution is then $a = 1,$ which gives us $b = -2,$ and $(a, b) = \boxed{\text{(C) } (1, -2)}$

 



\textbf{Solution 4 (Using the Answer Choices)}

Note that for roots $a$ and $b$, $ab = b$. This implies that $a=1$, and there is only one answer choice with $1$ in the position for $a$, hence, $(a,b) = \boxed{\mathrm{(C)}\ (1,-2)}$

 



\textbf{Solution 5(Answer Choices)}

We can use Vieta's formulas to solve this. We know that $a + b = -a$, and $ab = b.$ We can test the answer choices, so that we find that $\boxed{\mathrm{(C)}\ (1,-2)}$ works.

 -monkey\_land

 


	\item (\textit{AMC 10 10B 2002}) The product of three consecutive positive integers is $8$ times their sum. What is the sum of their squares?


 $\mathrm{(A) \ } 50\qquad \mathrm{(B) \ } 77\qquad \mathrm{(C) \ } 110\qquad \mathrm{(D) \ } 149\qquad \mathrm{(E) \ } 194$


 
 \textbf{Answer: }B



\textbf{Solution 1}

Let the three consecutive positive integers be $a-1$, $a$, and $a+1$. Since the mean is $a$, the sum of the integers is $3a$. So $8$ times the sum is just $24a$. With this, we now know that $a(a-1)(a+1)=24a\Rightarrow(a-1)(a+1)=24$.  $24=4\times6$, so $a=5$. Hence, the sum of the squares is $4^2+5^2+6^2=\boxed{\mathrm{ (B)}\ 77}$.

 


	\item (\textit{AMC 10 10B 2002}) For which of the following values of $k$ does the equation $\frac{x-1}{x-2} = \frac{x-k}{x-6}$ have no solution for $x$?


 $\textbf{(A) } 1\qquad \textbf{(B) } 2\qquad \textbf{(C) } 3\qquad \textbf{(D) } 4\qquad \textbf{(E) } 5$


 
 \textbf{Answer: }E



\textbf{Solution}

The domain over which we solve the equation is $\mathbb{R} \setminus \{2,6\}$.

 We can now cross-multiply to get rid of the fractions, we get $(x-1)(x-6)=(x-k)(x-2)$.

 Simplifying that, we get $7x-6 = (k+2)x - 2k$. Clearly for $k=5$ we get the equation $-6=-10$ which is never true. The answer is $\boxed{\mathrm{ (E)}\ 5}$

 
For other $k$, one can solve for $x$: $x(5-k) = 6-2k$, hence $x=\frac {6-2k}{5-k}$. We can easily verify that for none of the other 4 possible values of $k$ is this equal to $2$ or $6$, hence there is a solution for $x$ in each of the other cases.

 -Edited by XxHalo711 (typo within the solution)

 


	\item (\textit{AMC 10 10B 2002}) Find the value(s) of $x$ such that $8xy - 12y + 2x - 3 = 0$ is true for all values of $y$.


 $\textbf{(A) } \frac23 \qquad \textbf{(B) } \frac32 \text{ or } -\frac14 \qquad \textbf{(C) } -\frac23 \text{ or } -\frac14 \qquad \textbf{(D) } \frac32 \qquad \textbf{(E) } -\frac32 \text{ or } -\frac14$


 
 \textbf{Answer: }D



\textbf{Solution}

We have $8xy - 12y + 2x - 3 = 4y(2x - 3) + (2x - 3) = (4y + 1)(2x - 3)$. 

 As $(4y + 1)(2x - 3) = 0$ must be true for all $y$, we must have $2x - 3 = 0$, hence $\boxed{x = \frac 32\ \mathrm{ (D)}}$.

 



\textbf{Solution 2}

Since we want only the $y$-variable to be present, we move the terms only with the $y$-variable to one side, thus constructing $8xy - 12y + 2x - 3 = 0$ to $8xy + 2x = 12y + 3$. For there to be infinite solutions for $y$ and there is no $x$, we simply find a value of $x$ such that the equation is symmetrical. Therefore,

 $2x = 3$

 $8xy = 12y$

 There is only one solution, namely $x = \boxed{\dfrac{3}{2}}$ or $\text{(D)}$

 


	\item (\textit{AMC 10 10B 2002}) The number $25^{64}\cdot 64^{25}$ is the square of a positive integer $N$. In decimal representation, the sum of the digits of $N$ is


 $\mathrm{(A) \ } 7\qquad \mathrm{(B) \ } 14\qquad \mathrm{(C) \ } 21\qquad \mathrm{(D) \ } 28\qquad \mathrm{(E) \ } 35$


 
 \textbf{Answer: }B



\textbf{Solution}

Taking the root, we get $N=\sqrt{25^{64}\cdot 64^{25}}=5^{64}\cdot 8^{25}$.

 Now, we have $N=5^{64}\cdot 8^{25}=5^{64}\cdot (2^{3})^{25}=5^{64}\cdot 2^{75}$.

 Combining the $2$'s and $5$'s gives us $(2\cdot 5)^{64}\cdot 2^{(75-64)}=(2\cdot 5)^{64}\cdot 2^{11}=10^{64}\cdot 2^{11}$. 

 This is the number $2048$ with a string of sixty-four $0$'s at the end. Thus, the sum of the digits of $N$ is $2+4+8=14\Longrightarrow\boxed{\mathrm{ (B)}\ 14}$

 


	\item (\textit{AMC 10 10B 2002}) The positive integers $A$, $B$, $A-B$, and $A+B$ are all prime numbers. The sum of these four primes is


 $\mathrm{(A) \ } \text{even}\qquad \mathrm{(B) \ } \text{divisible by }3\qquad \mathrm{(C) \ } \text{divisible by }5\qquad \mathrm{(D) \ } \text{divisible by }7\qquad \mathrm{(E) \ } \text{prime}$


 
 \textbf{Answer: }E



\textbf{Solution 1}

Since $A-B$ and $A+B$ must have the same \href{https://artofproblemsolving.com/wiki/index.php?title=Parity}{parity}, and since there is only one even prime number, it follows that $A-B$ and $A+B$ are both odd. Thus one of $A, B$ is odd and the other even. Since $A+B > A > A-B > 2$, it follows that $A$ (as a prime greater than $2$) is odd. Thus $B = 2$, and $A-2, A, A+2$ are consecutive odd primes. At least one of $A-2, A, A+2$ is divisible by $3$, from which it follows that $A-2 = 3$ and $A = 5$. The sum of these numbers is thus $17$, a prime, so the answer is $\boxed{\mathrm{(E)}\ \text{prime}}$.

 



\textbf{Solution 2}

In order for both $A - B$ and $A + B$ to be prime, one of $A, B$ must be 2, or else both $A - B$, $A + B$ would be even numbers. 

 If $A = 2$, then $A < B$ and $A - B < 0$, which is not possible. Thus $B = 2$.

 Since $A$ is prime and $A > A - B > 2$, we can infer that $A > 3$ and thus $A$ can be expressed as $6n \pm 1$ for some natural number $n$. 

 However in either case, one of $A - B$ and $A + B$ can be expressed as $6n \pm 3 = 3(2n \pm 1)$ which is a multiple of 3. Therefore the only possibility that works is when $A - B = 3$ and 
 \[A + B + (A - B) + (A + B) = 5 + 2 + 3 + 7 = 17\]

 Which is a prime number. $\boxed{(E)}$

 



\textbf{Solution 3 (intuitive)}

Trying out some primes for $A$ and $B$ such that $A-B$ and $A+B$ are prime, $A=5$ and $B=2$ can be found almost immediately. Summing the four primes, the result is $5+2+3+7=17$, which is $\boxed{\mathrm{(E)}\ \text{prime}}$.

 


	\item (\textit{AMC 10 10B 2002}) For how many integers $n$ is $\frac{n}{20-n}$ the square of an integer?


 $\textbf{(A) } 1\qquad \textbf{(B) } 2\qquad \textbf{(C) } 3\qquad \textbf{(D) } 4\qquad \textbf{(E) } 10$


 
 \textbf{Answer: }D



\textbf{Solution 1}

Let $x^2 = \frac{n}{20-n}$, with $x \ge 0$ (note that the solutions $x < 0$ do not give any additional solutions for $n$). Then rewriting, $n = \frac{20x^2}{x^2 + 1}$. Since $\text{gcd}(x^2, x^2 + 1) = 1$, it follows that $x^2 + 1$ divides $20$. Listing the factors of $20$, we find that $x = 0, 1, 2 , 3$ are the only $\boxed{\mathrm{(D)}\ 4}$ solutions (respectively yielding $n = 0, 10, 16, 18$).

 



\textbf{Solution 2}

For $n<0$ and $n>20$ the fraction is negative, for $n=20$ it is not defined, and for $n\in\{1,\dots,9\}$ it is between 0 and 1.

 Thus we only need to examine $n=0$ and $n\in\{10,\dots,19\}$.

 For $n=0$ and $n=10$ we obviously get the squares $0$ and $1$ respectively.

 For prime $n$ the fraction will not be an integer, as the denominator will not contain the prime in the numerator.

 This leaves $n\in\{12,14,15,16,18\}$, and a quick substitution shows that out of these only $n=16$ and $n=18$ yield a square. Therefore, there are only $\boxed{\mathrm{(D)}\ 4}$ solutions (respectively yielding $n = 0, 10, 16, 18$).

 



\textbf{Solution 3}

If $\frac{n}{20-n} = k^2 \ge 0$, then $n \ge 0$ and $20-n > 0$, otherwise $\frac{n}{20-n}$ will be negative. Thus $0 \le n \le 19$ and 
 \[0 = \frac{0}{20-(0)} \le \frac{n}{20-n} \le \frac{19}{20-(19)} = 19\] Checking all $k$ for which $0 \le k^2 \le 19$, we have $0$, $1$, $2$, $3$ as the possibilities. $\boxed{(D)}$

 



\textbf{Solution 4}

For all integers x, $x^2$ is always a positive integer. So solve for $\frac{n}{20-n} = 0$, getting $n=0$ and $\frac{n}{20-n} = 1$, getting $n =10$. For all values n less than 0 and greater than 20, the value $\frac{n}{20-n}$ is negative, so now try values of n between 10 and 20. Quick substitution finds $0$, $10$, $16$, and $18$ which yields $x=0$, $x=1$, $x=2$, and $x=3$ respectively. 4 values, or $\boxed{(D)}$

 


 



\textbf{Solution 5}

Simon's Favourite Factoring Trick.

 Since $\frac{n}{20-n}$ is an integer $k$, we multiply both sides by $20-n$. This gives us $n=20k^2-nk^2$. We subtract $20k^2$ on both sides, then add $nk^2$ on both sides as a prerequisite for using Simon's Favorite Factoring Trick. We have $(k^2+1)(n-20)=20$. We then consider the different factors of $20$ that $k^2+1$ can be. It could be $1,2,4,5,10$, and $20$. After checking case by case, we then are able to identify that there are $4$ such $k$ values that also yield an integer $n$ value, meaning that there are $4$ values, so the correct answer is  $\boxed{(D)}$~CharmaineMa0729201

 



\textbf{Solution 6}

Similar to previous solutions, we can bound $n$ between $0$ and $20$. Then, we let $\frac{n}{20-n} = k^2$ for some integer $k$. Multiplying out the fraction and factoring, we get $n(k^2 + 1) = k^2 \cdot 20$. As before, $\text{gcd}(k^2 + 1, k^2) = 1$, so we have that $k^2 | n$. This means that $0 \leq k^2 \leq 20$, as $k^2$ now has to be at most $n$, so trying $k^2 = 0, 1, 4, 9, 16$, we get four possible $n$, namely $0$, $10$, $16$, and $18$. Thus, $\boxed{(D)}$.

 


	\item (\textit{AMC 10 10B 2002}) A regular octagon $ABCDEFGH$ has sides of length two.  Find the area of $\triangle ADG$.


 $\textbf{(A) } 4 + 2\sqrt2 \qquad \textbf{(B) } 6 + \sqrt2\qquad \textbf{(C) } 4 + 3\sqrt2 \qquad \textbf{(D) } 3 + 4\sqrt2 \qquad \textbf{(E) } 8 + \sqrt2$


 
 \textbf{Answer: }C



\textbf{Solution}


                    \begin{figure}[H]
                        \centering
                        \includegraphics[width=0.5\textwidth]{Solution/AMC_10_10B_2002_17_Solution_Figure_1.png}
                    \end{figure}
                    

 The area of the triangle $ADG$ can be computed as $\frac{DG \cdot AP}2$. We will now find $DG$ and $AP$.

 Clearly, $PFG$ is a right isosceles triangle with hypotenuse of length $2$, hence $PG=\sqrt 2$. The same holds for triangle $QED$ and its leg $QD$. The length of $PQ$ is equal to $FE=2$.Hence $GD = 2 + 2\sqrt 2$, and $AP = PD = 2 + \sqrt 2$.

 Then the area of $ADG$ equals $\frac{DG \cdot AP}2 = \frac{(2+2\sqrt 2)(2+\sqrt 2)}2 = \frac{8+6\sqrt 2}2 = \boxed{\textbf{(C)}=4+3\sqrt 2}$.

 



\textbf{Solution 2 (a bit of trigonometry)}

(Refer to diagram above)

 Notice how the area of triangle is equivalent to the area of the octagon minus $2A_{ABCD}$ and then minus $A_{\triangle AHG}$.

 THe trapezoid $ABCD$ is an isosceles trapezoid. It can be broken apart into a rectangle, and two right triangles. These triangles are isosceles, as $\angle GFP = \angle GFE - \angle PFE = 135^{\circ} - 90^{\circ} = 45^{\circ}$. 

 $\sin(45^{\circ}) = \frac{\sqrt{2}}{2}$, so the legs of the isosceles right triangles are $\sqrt{2}$. The area of the trapezoid can be calculated as followed:

 $A_{PQEF} + 2A_{GFP} = 2 \cdot \sqrt{2} + 2 \cdot \sqrt{2}^2 \div 2 = 2\sqrt{2}+2$.

 The area formula for a triangle is $\frac{1}{2}ab\sin(C)$, where $C$ is the included angle between sides $a$ and $b$. Using this for triangle $AHG$, we get $\frac{1}{2} \cdot 4 \cdot \sin(135)$. $\sin(135) = \sin(90+45) = \sin(90)\cos(45) + \cos(90)\sin(45) = 1 \cdot \frac{\sqrt{2}}{2} + 0 \cdot \frac{\sqrt{2}}{2} = \frac{\sqrt{2}}{2}$. Therefore, the area of triangle $AHG$ is $\sqrt{2}$.

 The total area between two times the trapezoid $ABCD$ and triangle $AHG$ is $5\sqrt{2}+4$. The area of the octagon is $2(1+\sqrt{2}) \cdot 2^2 = 8 + 8\sqrt{2}$. Subtracting the areas, we get $4+3\sqrt{2}$, or $\boxed{\text{C}}$.

 



\textbf{Solution 3 (trig bash)}


                    \begin{figure}[H]
                        \centering
                        \includegraphics[width=0.5\textwidth]{Solution/AMC_10_10B_2002_17_Solution_Figure_2.png}
                    \end{figure}
                    We are trying to find $\triangle ADG$ by using the formula $\frac{1}{2}ab\sin(C)$ and using sides $AG$ and $GD$.

 We note that the angle HGM is $\frac{45}{2}^{\circ}$. We need to find the length of $AG$ Then we recall the half angle identity:

 $\cos\left(\frac{\theta}{2}\right) = \sqrt{\frac{1+\cos(\theta)}{2}}$

 Then by substituting $\cos(45^\circ) = \frac{1}{\sqrt{2}}$ into the identity above, we get 

 $\cos\left(\frac{\cos(45/2)^\circ}{2}\right) = \sqrt{\frac{1+\cos(45)}{2}} = \sqrt{\frac{1+\frac{1}{\sqrt{2}}}{2}} = \sqrt{\frac{\sqrt{2}+1}{2\sqrt{2}}} = \sqrt{\frac{2+\sqrt{2}}{4}}$

 The length of $AG$ will be $2 \cdot 2 \cdot \sqrt{\frac{2+\sqrt{2}}{4}} = 4 \sqrt{\frac{2+\sqrt{2}}{4}}$. We will leave this without simplifying. 

 Now we find the length of GD.

 In the diagram above, $\triangle IFG$ and $\triangle JED$ are isosceles right triangles. Therefore, $IF = EJ = \frac{2}{\sqrt{2}} = \sqrt{2}$, so $GD = IJ = IF + FE + EJ = \sqrt{2} + 2 + \sqrt{2} = 2 + 2\cdot \sqrt{2}$

 Now we have the lengths of $AG$ and $GD$, we still have one step left: finding $\sin(\angle AGD)$. $\angle AGD =(90 - \angle HGM)^{\circ}$, because $\angle DGH = 90^\circ$. We already found $\angle HGM$, which is $\frac{45}{2}^{\circ}$. Therefore, $\angle AGD =(90 - \frac{45}{2})^{\circ}$. Now we can find the area of $\triangle ADG$ 

 Substituting our values into the formula $\frac{1}{2}ab\sin(C)$, we get

 $\text{Area of } \triangle ADG = \frac{1}{2} \cdot \sin\left((90-\frac{45}{2})^{\circ}\right) \cdot 4\sqrt{\frac{2+\sqrt{2}}{4}} \cdot (2 + 2\sqrt{2})$ 

 Using the identity $\sin(90^\circ - x) = \cos(x), \quad \sin\left((90-\frac{45}{2})^{\circ}\right) = \cos\left(\frac{45}{2}^\circ\right)$. 

 Substituting this back into the expression $\frac{1}{2} \cdot \sin\left((90-\frac{45}{2})^{\circ}\right) \cdot 4\sqrt{\frac{2+\sqrt{2}}{4}} \cdot (2 + 2\sqrt{2})$ , we get 

 $\frac{1}{2} \cdot \cos\left(\frac{45}{2}^\circ\right) \cdot 4\sqrt{\frac{2+\sqrt{2}}{4}} \cdot (2 + 2\sqrt{2}) = \frac{1}{2} \cdot \sqrt{\frac{2+\sqrt{2}}{4}} \cdot 4\sqrt{\frac{2+\sqrt{2}}{4}} \cdot (2 + 2\sqrt{2})=$$2 \cdot \frac{2+\sqrt{2}}{4} \cdot (2 + 2\sqrt{2}) = \frac{(2+\sqrt{2})(2+2\sqrt{2})}{2} =  \frac{4 + 4\sqrt{2} + 2\sqrt{2} + 4}{2} = \frac{8 + 6\sqrt{2}}{2} = 4 + 3\cdot \sqrt{2}$

 Therefore, the answer is $\boxed{(C)}$

 


	\item (\textit{AMC 10 10B 2002}) Four distinct circles are drawn in a plane. What is the maximum number of points where at least two of the circles intersect?


 $\textbf{(A) } 8\qquad \textbf{(B) } 9\qquad \textbf{(C) } 10\qquad \textbf{(D) } 12\qquad \textbf{(E) } 16$


 
 \textbf{Answer: }D



\textbf{Solution 1}

For any given pair of circles, they can intersect at most $2$ times. Since there are ${4\choose 2} = 6$ pairs of circles, the maximum number of possible intersections is $6 \cdot 2 = 12$. We can construct such a situation as below, so the answer is $\boxed{\mathrm{(D)}\ 12}$.

 
                    \begin{figure}[H]
                        \centering
                        \includegraphics[width=0.5\textwidth]{Solution/AMC_10_10B_2002_18_Solution_Figure_1.png}
                    \end{figure}
                    

 



\textbf{Solution 2}

Because a pair or circles can intersect at most $2$ times, the first circle can intersect the second at $2$ points, the third can intersect the first two at $4$ points, and the fourth can intersect the first three at $6$ points. This means that our answer is $2+4+6=\boxed{\mathrm{(D)}\ 12}.$

 



\textbf{Solution 3}

Pick a circle any circle- $4$ ways. Then, pick any other circle- $3$ ways. For each of these circles, there will be $2$ intersections for a total of $4*3*2$ = $24$ intersections. However, we have counted each intersection twice, so we divide for overcounting. Therefore, we reach a total of $\frac{24}{2}=\boxed{12}$, which corresponds to $\text{(D)}$.

 


	\item (\textit{AMC 10 10B 2002}) Suppose that $\{a_n\}$ is an arithmetic sequence with

 \[a_1+a_2+\cdots+a_{100}=100 \text{ and } a_{101}+a_{102}+\cdots+a_{200}=200.\]
What is the value of $a_2 - a_1 ?$


 $\mathrm{(A) \ } 0.0001\qquad \mathrm{(B) \ } 0.001\qquad \mathrm{(C) \ } 0.01\qquad \mathrm{(D) \ } 0.1\qquad \mathrm{(E) \ } 1$


 
 \textbf{Answer: }C



\textbf{Solution 1(Most efficient)}

We should realize that the two equations are 100 terms apart, so by subtracting the two equations in a form like...

 
 \[(a_{101} - a_1) + (a_{102} - a_2) +... + (a_{200} - a_{100}) = 200-100 = 100\]

 ...we get the value of the common difference of every hundred terms one hundred times. So we have to divide the answer by one hundred to get ... 

 $\frac{100}{100} = 1$

 ...the common difference of every hundred terms. Then we have to simply divide the answer by hundred again to find the common difference between one term, therefore...

 $\frac{1}{100} =\boxed{(\textbf{C})\ 0.01}$

 



\textbf{Solution 2}

Adding the two given equations together gives 

 $a_1+a_2+...+a_{200}=300$.

 Now, let the common difference be $d$. Notice that $a_2-a_1=d$, so we merely need to find $d$ to get the answer. The formula for an arithmetic sum is 

 $\frac{n}{2}\left(2a_1+d(n-1)\right)$,

 where $a_1$ is the first term, $n$ is the number of terms, and $d$ is the common difference. Now we use this formula to find a closed form for the first given equation and the sum of the given equations. For the first equation, we have $n=100$. Therefore, we have 

 $50(2a_1+99d)=100$,

 or 

 $2a_1+99d=2$.          *\textbf{(1)}

 For the sum of the equations (shown at the beginning of the solution) we have $n=200$, so

 $100(2a_1+199d)=300$

 or 

 $2a_1+199d=3$          *\textbf{(2)}

 Now we have a system of equations in terms of $a_1$ and $d$.Subtracting \textbf{(1)} from \textbf{(2)} eliminates $a_1$, yielding $100d=1$, and $d=a_2-a_1=\frac{1}{100}=\boxed{(\textbf{C})\ 0.01}$.

 



\textbf{Solution 3}

Subtracting the 2 given equations yields

 
$(a_{101}-a_1)+(a_{102}-a_2)+(a_{103}-a_3)+...+(a_{200}-a_{100})=100$

 
Now express each $a_n$ in terms of first term $a_1$ and common difference $x$ between consecutive terms

 
$((a_1+100x)-(a_1))+((a_1+101x)-(a_1+x))+((a_1+102x)-(a_1+2x))+...+((a_1+199x)-(a_1+99x))=100$

 
Simplifying and canceling $a_1$ and $x$ terms gives

 
$100x+100x+100x+...+100x=100$

 
$100x\times100=100$ 

 
$100x=1$

 
$x=0.01=\boxed{(\textbf{C})\ 0.01}$

 



\textbf{Solution 4}

Use the Summation Formula, $S = \frac{n}{2}(2a+(n-1)d)$. For the first equation given, you can plug in values as such: 
 \[100 = \frac{100}{2}(2a+99d).\] Solving this gets $2a+99d = 2$. Doing the same for the second equation given, you get 
 \[200 = \frac{100}{2}(2a+99d).\] or $2a+99d=4$. Well, notice that you started summing from $a_{101}$ rather than $a_1$. So, use the explicit formula to get $a_{101} = a_1 + (101-1)d$, or $a_{101} = a + 100d$. Substitution gets $2a + 299d = 4$. Now, we need the value of $d$, as $a_2 - a_1 = d$. Use elimination on the two equations found to get $200d = 2$. Therefore, $d = \frac{1}{100}$, or $\boxed{(\textbf{C})\ 0.01}$.

 -N.

 


	\item (\textit{AMC 10 10B 2002}) Let $a, b,$ and $c$ be real numbers such that $a-7b+8c=4$ and $8a+4b-c=7.$ Then $a^2-b^2+c^2$ is


 $\mathrm{(A) \ } 0\qquad \mathrm{(B) \ } 1\qquad \mathrm{(C) \ } 4\qquad \mathrm{(D) \ } 7\qquad \mathrm{(E) \ } 8$


 
 \textbf{Answer: }B



\textbf{Solution 1}

Rearranging, we get $a+8c=7b+4$ and $8a-c=7-4b$

 Squaring both, $a^2+16ac+64c^2=49b^2+56b+16$ and $64a^2-16ac+c^2=16b^2-56b+49$ are obtained.

 Adding the two equations and dividing by $65$ gives $a^2+c^2=b^2+1$, so $a^2-b^2+c^2=\boxed{(\text{B})1}$.

 



\textbf{Solution 2}

The easiest way is to assume a value for $a$ and then solve the system of equations. For $a = 1$, we get the equations $-7b + 8c = 3$ and$4b - c = -1$Multiplying the second equation by $8$, we have $32b - 8c = -8$Adding up the two equations yields $25b = -5$, so $b = -\frac{1}{5}$We obtain $c = \frac{1}{5}$ after plugging in the value for $b$.Therefore, $a^2-b^2+c^2 = 1-\frac{1}{25}+\frac{1}{25}=\boxed{1}$ which corresponds to $\text{(B)}$.This time-saving trick works only because we know that for any value of $a$, $a^2-b^2+c^2$ will always be constant (it's a contest), so any value of $a$ will work. This is also called \href{https://artofproblemsolving.com/wiki/index.php?title=Without\_loss\_of\_generality}{without loss of generality} or WLOG.

 



\textbf{Solution 3 (fake solve)}

Notice that the coefficients of $a$ and $c$ are pretty similar (15s for reading and noticing), so let $b=0$ gives $a+8c=4$, and $8a-c=7$ (10s writing). Since the desired quantity simplifies to $a^2+c^2$, the $ac$ term of the quadratics after squaring gets canceled by adding up the squares of the two equations because they have the same coefficients but opposite sign (15s mind-binom). This simplifies to $65(a^2+c^2)=16+49$, or $a^2-b^2+c^2=\frac{65}{65}=\boxed{(\text{B})1}$(15s writing and addition and fraction simplification and (B) circling and submission)

 



\textbf{Solution 4 (very simple)}

Note that the answer choices list different numbers, as opposed to an expression involving variables. This means that the expression $a^2-b^2+c^2$ evaluates to a constant. This means that we can assume $a = 0$ since the expression will stay the same for any value of a. Then since $a = 0$ and $8a = 8 \cdot 0 = 0$, the equations given in the problem become $-7b + 8c = 4$ and $4b - c = 7$. Multiplying the second equation, we obtain $32b - 8c = 56$. Adding this to our first equation, this gives us $25b = 60$, so $b = 60/25 = 12/5$. Plugging this into $4b - c = 7$, we get $4(12/5) - c = 7 \implies 48/5 - c = 35/5 \implies c = 13/5$. Then, $a^2-b^2+c^2 = 0^2 - (12/5)^2 + (13/5)^2 = -144/25 + 169/25 = 25/25 = \boxed{(\text{B})1}$.

 
- Loquacious\_Autodidact

 


	\item (\textit{AMC 10 10B 2002}) Andy's lawn has twice as much area as Beth's lawn and three times as much as Carlos' lawn. Carlos' lawn mower cuts half as fast as Beth's mower and one third as fast as Andy's mower. If they all start to mow their lawns at the same time, who will finish first?


 $\mathrm{(A) \ } \text{Andy}\qquad \mathrm{(B) \ } \text{Beth}\qquad \mathrm{(C) \ } \text{Carlos}\qquad \mathrm{(D) \ } \text{Andy and Carlos tie for first.}\qquad \mathrm{(E) \ } \text{All three tie.}$


 
 \textbf{Answer: }B



\textbf{Solution 1}

We say Andy's lawn has an area of $x$. Beth's lawn thus has an area of $\frac{x}{2}$, and Carlos's lawn has an area of $\frac{x}{3}$. 

 We say Andy's lawn mower cuts at a speed of $y$. Carlos's cuts at a speed of $\frac{y}{3}$, and Beth's cuts at a speed $\frac{2y}{3}$.

 Each person's lawn is cut at a time of $\frac{\text{area}}{\text{rate}}$, so Andy's is cut in $\frac{x}{y}$ time, as is Carlos's. Beth's is cut in $\frac{3}{4}\times\frac{x}{y}$, so the first one to finish is $\boxed{\mathrm{(B)}\ \text{Beth}}$.

 



\textbf{Solution 2}

WLOG, we can set values of their lawns' areas and their owners' speeds. Let the area of Andy's lawn be $6$ units, Beth's lawn be $3$ units, and Carlos's lawn be $2$ units. Let Carlos's mowing area per hour (honestly the time you set won't matter) be $1$, let Beth's mowing area per hour be $2$, and let Andy's mowing area per hour be $3$. Now, we can easily calculate their time  by dividing their lawns' areas by their respective owners' speeds. Andy's time is $\frac{6}{3} = 2$ hours, Beth's time is $\frac{3}{2} = 1.5$ hours, and Carlos's time is $\frac{6}{3} = 2$ hours. Our answer is clearly $\boxed{\mathrm{(B)}\ \text{Beth}}$.~Solution by virjoy2001

 



\textbf{Solution 3(convert into a racing problem)}

We first set the area for Carlo's lawn to $x$. That leads to Beth's lawn's area to be $3x$ and Andy's lawn $6x$. After that, we set the speed of Carlo's lawn mower to $y$. We then get Beth's speed in terms of $y$, which is $2y$. Andy's lawn mower's speed is $3y$. WLOG, we set the width of their lawns to $0$. We then can imagine that it's a racetrack, with the distance being the area of their lawns, and they are competing with different speeds and distances. We then proceed by applying the speed formula, distance = speed*time. Using WLOG again, we can set Carlo's speed to $10$ mph and we can assume Carlo's lawn, leaving Beth with $20$ mph/$3$miles and Andy with $30$ mph/$6$ miles. After plugging in the variables, we clearly see that $\boxed{\mathrm{(B)}\ \text{Beth}}$ will win.~Calvin67

 


	\item (\textit{AMC 10 10B 2002}) Let $\triangle{XOY}$ be a right-angled triangle with $\angle{XOY}=90^\circ$. Let $M$ and $N$ be the midpoints of the legs $OX$ and $OY$, respectively. Given $XN=19$ and $YM=22$, find $XY$.


 $\mathrm{(A) \ } 24\qquad \mathrm{(B) \ } 26\qquad \mathrm{(C) \ } 28\qquad \mathrm{(D) \ } 30\qquad \mathrm{(E) \ } 32$


 
 \textbf{Answer: }B



\textbf{Solution 1}


                    \begin{figure}[H]
                        \centering
                        \includegraphics[width=0.5\textwidth]{Solution/AMC_10_10B_2002_22_Solution_Figure_1.png}
                    \end{figure}
                    

 Let $OM = x$, $ON = y$. By the \href{https://artofproblemsolving.com/wiki/index.php?title=Pythagorean\_Theorem}{Pythagorean Theorem} on $\triangle XON, MOY$ respectively,
 \begin{align*} (2x)^2 + y^2 &= 19^2\\ x^2 + (2y)^2 &= 22^2\end{align*}


 Summing these gives $5x^2 + 5y^2 = 845 \Longrightarrow x^2 + y^2 = 169$. 

 By the Pythagorean Theorem again, we have 

 
 \[(2x)^2 + (2y)^2 = XY^2 \Longrightarrow XY = \sqrt{4(x^2 + y^2)} = \sqrt{4(169)} = \sqrt{676} = \boxed{\mathrm{(B)}\ 26}.\]

 Alternatively, we could note that since we found $x^2 + y^2 = 169$, segment $MN=13$. Right triangles $\triangle MON$ and $\triangle XOY$ are similar by Leg-Leg with a ratio of $\frac{1}{2}$, so $XY=2(MN)=\boxed{\mathrm{(B)}\ 26}$.

 


 



\textbf{Solution 2}

Let $XO=x$ and $YO=y.$ Then, $XY=\sqrt{x^2+y^2}.$ 

 Since $XN=19$ and $YM=22,$ 
 \[XN^2=19^2=x^2+(\dfrac{y}{2})^2)=\dfrac{x^2}{4}+y^2\]
 \[YM^2=22^2=(\dfrac{x}{2})^2+y^2=\dfrac{x^2}{4}+y^2.\]

 Adding these up: 
 \[19^2+22^2=\dfrac{4x^2+y^2}{4}+\dfrac{x^2+4y^2}{4}\]
 \[845=\dfrac{5x^2+5y^2}{4}\]
 \[3380=5x^2+5y^2\]
 \[676=x^2+y^2.\] 

 Then, we substitute: $XY=\sqrt{x^2+y^2}=\sqrt{676}=\boxed{26}.$

 


 



\textbf{Solution 3 (Solution 1 but shorter)}

Refer to the diagram in solution 1. $4x^2+y^2=361$ and $4y^2+x^2=484$, so add them: $5x^2+5y^2=845$ and divide by 5: $x^2+y^2=169$ so $\dfrac{XY}{2}=\sqrt{169}=13$ and so $XY=26$, or answer $\boxed{\text{(B)}}$.

 



\textbf{Solution 4}

Use the diagram in solution 1. Get $4x^2+y^2=361$ and $4y^2+x^2=484$, and multiply the second equation by 4 to get $4x^2+16y^2=1936$ and then subtract the first from the second. Get $15y^2=1575$ and $y^2=105$. Repeat for the other variable to get $15x^2=960$ and $x^2=64$. Now XY is equal to the square root of four times these quantities, so $(105+64) \cdot 4=676$, and $XY=\boxed{\text{(B)}\ 26}$.

 



\textbf{Solution 5 (Quick and good for in-contest use)}

Let $XM = MO = y$, and $ON = NY  = x$. Now, we can write some equations. We have that 
 \[(2y)^2+x^2 = 19^2 \Longrightarrow 4y^2 + x^2 = 361,\] and 
 \[(2x)^2+y^2 = 22^2 \Longrightarrow 4x^2+y^2 = 484.\] Adding the $2$ equations up we get 
 \[5x^2+5y^2 = 845.\] The hypotenuse of the triangle is equal to 
 \[\sqrt{4x^2+4y^2},\] so that is what we want. In order to get 
 \[4x^2+4y^2,\] we need to subtract $x^2+y^2$ from $5x^2+5y^2$, and since $x^2+y^2 = \dfrac{1}{5}\left(5x^2+5y^2\right)$, we have that 
 \[4x^2+4y^2 = \dfrac{4}{5} \cdot 845 = 676.\] Therefore, the hypotenuse is $\sqrt{676} = \boxed{\text{(B)}\ 26}$.

 -jb2015007

 


	\item (\textit{AMC 10 10B 2002}) Let $\{a_k\}$ be a sequence of integers such that $a_1=1$ and $a_{m+n}=a_m+a_n+mn,$ for all positive integers $m$ and $n.$ Then $a_{12}$ is


 $\mathrm{(A) \ } 45\qquad \mathrm{(B) \ } 56\qquad \mathrm{(C) \ } 67\qquad \mathrm{(D) \ } 78\qquad \mathrm{(E) \ } 89$


 
 \textbf{Answer: }D



\textbf{Solution 1}

When $m=1$, $a_{n+1}=1+a_n+n$. Hence,
 \[a_{2}=1+a_1+1\]
 \[a_{3}=1+a_2+2\]
 \[a_{4}=1+a_3+3\]
 \[\dots\]
 \[a_{12}=1+a_{11}+11\]Adding these equations up, we have that $a_{12}=12+(1+2+3+...+11)=\boxed{\mathrm{(D) \ } 78}$

 



\textbf{Solution 2}

Substituting $n=1$ into $a_{m+n}=a_m+a_n+mn$: $a_{m+1}=a_m+a_{1}+m$.

 Since $a_1 = 1$, $a_{m+1}=a_m+m+1$.

 Therefore, $a_m = a_{m-1} + m, a_{m-1}=a_{m-2}+(m-1), a_{m-2} = a_{m-3} + (m-2)$, and so on until $a_2 = a_1 + 2$.

 Adding the Left Hand Sides of all of these equations gives $a_m + a_{m-1} + a_{m-2} + a_{m-3} + \cdots + a_2$.

 Adding the Right Hand Sides of these equations gives 

 $(a_{m-1} + a_{m-2} + a_{m-3} + \cdots + a_1) + (m + (m-1) + (m-2) + \cdots + 2)$.

 These two expressions must be equal; hence $a_m + a_{m-1} + a_{m-2} + a_{m-3} + \cdots + a_2 = (a_{m-1} + a_{m-2} + a_{m-3} + \cdots + a_1) + (m + (m-1) + (m-2) + \cdots + 2)$ and $a_m = a_1 + (m + (m-1) + (m-2) + \cdots + 2)$.

 Substituting $a_1 = 1$: $a_m = 1 + (m + (m-1) + (m-2) + \cdots + 2) = 1+2+3+4+ \cdots +m = \frac{(m+1)(m)}{2}$.

 Thus we have a general formula for $a_m$ and substituting $m=12$: $a_{12} = \frac{(13)(12)}{2} = (13)(6) = \boxed{\mathrm{(D)} 78}$.

 



\textbf{Solution 3}

We can literally just plug stuff in. No prerequisite is actually said in the sequence. Since $a_{m+n} = a_m+a_n +mn$, we know $a_2=a_1+a_1+1\cdot1=1+1+1=3$. After this, we can use $a_2$ to find $a_4$. $a_4=a_2+a_2+2\cdot 2 = 3+3+4 = 10$. Now, we can use $a_2$ and $a_4$ to find $a_6$, or $a_6=a_4+a_2+4\cdot 2 = 10+3+8=21$. Lastly, we can use $a_6$ to find $a_{12}$. $a_{12} = a_6+a_6+6\cdot 6 = 21+21+36= \boxed{\mathrm{(D) \ } 78}$

 



\textbf{Solution 4}

We can set $n$ equal to $m$, so we can say that 
 \[a_{m + m} = a_m + a_m + m*m\]
 \[a_{2m} = 2a_m + m^2\]

 We set $2m = 12$, we get $m = 6$.
 \[a_{12} = 2a_6 + 36\]

 We set $2m = 6$m, we get $m = 3$.
 \[a_6 = 2a_3 + 9\]

 Solving for $a_3$ is easy, just direct substitution.
 \[a_2 = 1 + 1 + 1 = 3\]
 \[a_3 = a_{2 + 1} = 3 + 1 + 2 = 6\]

 Substituting, we get
 \[a_6 = 2(6) + 9 = 21\]
 \[a_{12} = 2(21) + 36 = 78\]

 Thus, the answer is $\boxed{D}$.

 



\textbf{Solution 5}

Note that the sequence of triangular numbers $T_n=1+2+3+...+n$ satisfies these conditions. It is immediately obvious that it satisfies $a_1=1$, and $a_{m+n}=a_m+a_n+mn$ can be visually proven with the diagram below.

 
                    \begin{figure}[H]
                        \centering
                        \includegraphics[width=0.5\textwidth]{Solution/AMC_10_10B_2002_23_Solution_Figure_1.png}
                    \end{figure}
                    

 This means that we can use the triangular number formula $T_n = \frac{n(n+1)}{2}$, so the answer is $T_{12} = \frac{12(12+1)}{2} = \boxed{\mathrm{(D) \ } 78}$.

 



\textbf{Solution 6}

Follow solution 4 with plugging in m=n, but for the last part since since the subscript is 2n, and we aren't given anything that has odd numbers ( without substitution into the original equation), we need to find two numbers that add up to 12, which result from the multiplication of two even numbers. We realize that 4+8 works, since 4=2x2 and 8=4x2. We then proceed in substituting n=2 and n=4 into the equation, then add to obtain 78.~Charmainema07292010

 



\textbf{Solution 6 (Fast - Good for use in contest)}

Find $a_2$ which is just 
 \[a_2 = a_{1+1} = 1+1+1 = 3\]. Now, we can find $a_4$ and so on until $a_{12}$. So here goes: 
 \[a_4 = a_{2+2} = 3+3+4 = 10\]
 \[a_8 = a_{4+4} = 10+10+16 = 36\]
 \[a_{12} = a_{8+4} = 36+10+32 = \boxed{\mathrm{(D) \ } 78}\]

 -jb2015007

 


	\item (\textit{AMC 10 10B 2002}) Riders on a Ferris wheel travel in a circle in a vertical plane. A particular wheel has radius $20$ feet and revolves at the constant rate of one revolution per minute. How many seconds does it take a rider to travel from the bottom of the wheel to a point $10$ vertical feet above the bottom?


 $\mathrm{(A) \ } 5\qquad \mathrm{(B) \ } 6\qquad \mathrm{(C) \ } 7.5\qquad \mathrm{(D) \ } 10\qquad \mathrm{(E) \ } 15$


 
 \textbf{Answer: }D



\textbf{Solution 1}


                    \begin{figure}[H]
                        \centering
                        \includegraphics[width=0.5\textwidth]{Solution/AMC_10_10B_2002_24_Solution_Figure_1.png}
                    \end{figure}
                    

 We can let this circle represent the ferris wheel with center $O,$ and $C$ represent the desired point $10$ feet above the bottom. Draw a diagram like the one above. We find out $\triangle OBC$ is a $30-60-90$ triangle. That means $\angle BOC = 60^\circ$ and the ferris wheel has made $\frac{60}{360} = \frac{1}{6}$ of a revolution. Therefore, the time it takes to travel that much of a distance is $\frac{1}{6}\text{th}$ of a minute, or $10$ seconds. The answer is $\boxed{\mathrm{(D) \ } 10}$. Alternatively, we could also say that $\triangle ABC$ is congruent to $\triangle OBC$ by SAS, so $AC$ is 20, and $\triangle AOC$ is equilateral, and $\angle BOC = 60^\circ$

 



\textbf{Solution 2 (For trigonometry students only)}

The path that the rider takes along the Ferris wheel can be represented by a sinusoidal graph, where $x$ represents the time in seconds. Since $x=0$ is at the crest of the graph and not at the midline, we will use a cosine graph. Therefore, we will use the form:
 \[f(x) = A\cos(Bx - C) + D.\]

 The graph starts at the lowest point at $0$ feet, then goes up to reach the highest point at $40$ feet, then comes back down. Therefore, the amplitude is $20$ (and negative since it starts at the bottom, not the top), and the vertical shift is $20$. There is no horizontal shift since the lowest point is at $x=0$. It takes $60$ seconds to make one full revolution (the period), so $B = \frac{2\pi}{60} = \frac{\pi}{30}.$ Now we have all the parts we need for the equation of our graph, and we can set it equal to the height we want, $10.$

 
 \[10 = -20\cos\left(\frac{\pi}{30}x\right) + 20.\]

 We get to $\frac{1}{2} = \cos\left(\frac{\pi}{30}x\right)$ and remember that the problem is looking for the first instance of $f(x) = 10$, so $\frac{\pi}{30}x = \frac{\pi}{3}.$ Solving, we get that $x = \boxed{\mathrm{(D) \ } 10}$.

 


	\item (\textit{AMC 10 10B 2002}) When $15$ is appended to a list of integers, the mean is increased by $2$. When $1$ is appended to the enlarged list, the mean of the enlarged list is decreased by $1$. How many integers were in the original list?


 $\mathrm{(A) \ } 4\qquad \mathrm{(B) \ } 5\qquad \mathrm{(C) \ } 6\qquad \mathrm{(D) \ } 7\qquad \mathrm{(E) \ } 8$


 
 \textbf{Answer: }A



\textbf{Solution 1}

Let $x$ be the sum of the integers and $y$ be the number of elements in the list. Then we get the equations $\dfrac{x+15}{y+1}=\dfrac{x}{y}+2$ and $\dfrac{x+15+1}{y+1+1}=\dfrac{x+16}{y+2}=\frac{x}{y}+2-1=\frac{x}{y}+1$. With a lot of algebra, the solution is found to be $y= \boxed{\textbf{(A)}\ 4}$.

 



\textbf{Solution 2}

We let $n$ be the original number of elements in the set and we let $m$ be the original average of the terms of the original list. Then we have $mn$ is the sum of all the elements of the list. So we have two equations: 
 \[mn+15=(m+2)(n+1)=mn+m+2n+2\] and 
 \[mn+16=(m+1)(n+2)=mn+2m+n+2.\]Simplifying both equations and we get,
 \[13=m+2n\]
 \[14=2m+n\]Solving for $m$ and $n$, we get $m=5$ and $n=\boxed{\textbf{(A)}4}$.

 



\textbf{Solution 3}

Warning: This solution will rarely ever work in any other case. However, seeing that you can so easily plug and chug in problem 25 it is funny to see this.

 Plug and chug random numbers with the answer choices, starting with the choice of $4$ numbers. You see that if you have 4 5s and you add 15 to the set, the resulting mean will be 7; we can verify this with math
 \[\frac{5+5+5+5+15}{5}=7\]adding in 1 to the set you result in the mean to be 6.
 \[\frac{5+5+5+5+15+1}{6}=6\]Thus we conclude that 4 is the correct choice or $\boxed{\textbf{(A)}}$

 


\end{enumerate}